% Root proposed complete B.116. Independent review required before integration.
\subsection*{An intrinsic obstruction in the actual regional group}

\statement{B.116}{Proved}{Two exact discrete regional centralizers exclude every automorphism extending the prefix change}
\label{stmt:B116}
Let $R$ be the actual annular bad-region group of B.112, before the
two cross relations of B.114. We retain its marked generators
$a,b,C,d,e,h,k$, put $B=CbC$, and set $z=Ch$.
The capital $C$ denotes B.112's adapted cone.
All groups and centralizers in this statement are discrete. Then
\begin{align}
 C_R(a)&=\langle a\rangle\times
   \bigl(\mathbb Z^2(b,k)*F(CaC,CdC,CeC)\bigr),
   \label{eq:B116-centralizer-a}\\
 C_R(aB)&=\langle aB\rangle\times F(b,CaC).
   \label{eq:B116-centralizer-prefix}
\end{align}
Both displayed cyclic factors are infinite. Their centralizer
abelianizations have ranks six and three, respectively. Consequently
no automorphism of $R$ sends $a$ to $aB$. In particular the actual
prefix automorphism $\theta$ of the common free neck $K$ in B.115
extends to no discrete automorphism of $R$, with no restriction on
the proposed image of $k$.

This statement does not identify centralizers of these infinite-order
elements in $\widehat R$. B.115 supplies the separate full profinite
obstruction for homomorphisms whose image of $k$ is ordinary and has
odd base winding. Neither statement transfers a regional nonmembership
claim through the two actual cross relations.

\begin{proof}
\emph{B.116.1. The entire regional graph kernel.}
Let
\[
 D=\langle C,k\mid C^2=1\rangle=C_2*\mathbb Z,
 \qquad T=\langle k\rangle.
\]
B.112's complete marked presentation, with $h=Cz$, has the generic
relations
\[
 C^2=1,\qquad [C,z]=[k,z]=1,\qquad
 [k,U]=1\quad(U=a,b,d,e),\qquad f=Czk^{-1}.
\]
Define a graph $\Omega$ on the vertices
\[
 (gT,U,n),\qquad gT\in D/T,\quad
 U\in\{a,b,d,e\},\quad n\in\mathbb Z.
\]
Its edges are all left $D$-translates of the following ten types,
at every integer $n$:
\begin{align*}
 &(T,U,n)\;--\;(CT,U,n) &&(U=a,b,d,e),\\
 &(T,a,n)\;--\;(T,b,n),\\
 &(CT,a,n)\;--\;(T,d,n),\qquad
   (CT,a,n)\;--\;(T,e,n),\\
 &(T,d,n)\;--\;(CT,e,n),\\
 &(T,b,n)\;--\;(CT,d,n+1),\qquad
   (T,b,n)\;--\;(CT,e,n+1).
\end{align*}
Write $A(\Omega)$ for the graph group. The actions of $D$ on cosets
and of $z$ on the integer coordinate commute and preserve every edge.
There is an exact presentation isomorphism
\begin{equation}\label{eq:B116-regional-kernel}
 R=A(\Omega)\rtimes(D\times\mathbb Z_z).
\end{equation}
Indeed, send the graph vertex $(gT,U,n)$ to
$gz^nUz^{-n}g^{-1}$. This is independent of the coset representative
because $k$ commutes with $z$ and with $U$. Every edge becomes a
conjugate of one of the ten actual node relations of B.112: the two
interlevel relations use $h=Cz$. Conversely, send $C,k,z$ to the
acting factor and each $U$ to $(T,U,0)$. The six generic relations,
the cone square, and all ten node relations hold. The two maps are
inverse on these generators and on every graph vertex. Thus
\eqref{eq:B116-regional-kernel} includes every coarse and height
translate, with no omitted kernel relations.

Put $N=A(\Omega)$. The actual elements $a$ and $B$ are the vertices
$(T,a,0)$ and $(CT,b,0)$. The ten edge types give their complete links:
\begin{align}
 \operatorname{lk}(a)
  &=\{(T,b,0)\}\ \cup\
    \{(k^mCT,U,0):m\in\mathbb Z,\ U=a,d,e\},
    \label{eq:B116-link-a}\\
 \operatorname{lk}(B)
  &=\{(CT,a,0)\}\ \cup\
    \{(Ck^mCT,b,0),(Ck^mCT,d,1),(Ck^mCT,e,1):m\in\mathbb Z\}.
    \label{eq:B116-link-B}
\end{align}
For example, a translated diagonal through $(T,a,0)$ has its other
endpoint $(k^mCT,a,0)$; the two same-height cross types give its
$d,e$ companions. This accounts for all edge types containing $a$.
For $B$, the diagonal gives the $b$ family, its same-coset edge gives
$(CT,a,0)$, and the two forward edges give the height-one $d,e$ families.

The free-product normal form in $C_2*\mathbb Z$ shows that the
cosets $k^mCT$ are distinct and that
\[
 Ck^mC\in T\quad\Longleftrightarrow\quad m=0.
\]
Consequently the full common link is
\begin{equation}\label{eq:B116-common-link}
 \operatorname{lk}(a)\cap\operatorname{lk}(B)
       =\{(T,b,0),(CT,a,0)\}.
\end{equation}
These two vertices are nonadjacent.

The entire induced graph on \eqref{eq:B116-link-a} is also discrete.
To see this for all $m$ at once, use the homomorphism $D\to C_2$
which sends $C$ to one and $k$ to zero. Every listed $a,d,e$ vertex
has odd coset parity. Every edge between these types has opposite
coset parities at its endpoints. The only same-coset edge uses
types $a,b$, whereas the sole $b$ vertex here has coset $T$.
Finally, the interlevel edges cannot occur among height-zero vertices.
This proves discreteness without truncating the coset or height set.

\emph{B.116.2. The needed centralizer identities for any graph.}
We prove the special case used here directly. Let $x,y$ be nonadjacent
vertices of any graph $\Gamma$, which need not be finite. Set
$H=A(\Gamma\setminus\{x\})$ and $L=A(\operatorname{lk}(x))$.
The graph presentation is the HNN extension
\[
 A(\Gamma)=\langle H,x\mid x\ell x^{-1}=\ell\ (\ell\in L)\rangle.
\]
These are genuine full-graph inclusions: deleting the other generators
gives explicit retractions. Its Bass--Serre normal form shows that
both $x$ and $xy$ have translation length one. A centralizer preserves
the corresponding axis and acts on it by translations, since an axis
reflection cannot commute with a nonzero translation. The indicated
length-one element generates the image of that action. Its kernel
fixes the axis pointwise and commutes with that element.

For $x$ the pointwise stabilizer is $L$. For $xy$, the first two edge
stabilizers force it into
$L\cap xyLy^{-1}x^{-1}=L\cap yLy^{-1}$, since $x$ centralizes $L$.
Now split the same graph group with stable letter $y$. The subgroup
$L$ is in its base group because $y\notin\operatorname{lk}(x)$.
Britton reduction says that $yvy^{-1}$, with $v\in L$, belongs to
that base only when $v\in A(\operatorname{lk}(y))$; then $y$ commutes
with $v$. Thus
\[
 L\cap yLy^{-1}
   =A(\operatorname{lk}(x)\cap\operatorname{lk}(y)).
\]
The intersection identity for the two full graph subgroups follows
from their commuting deletion retractions. Conversely this common-link
group commutes with both $x,y$ and fixes the initial edge, hence fixes
the entire axis. We have proved the exact formulas
\begin{align}
 C_{A(\Gamma)}(x)&=\langle x\rangle\times A(\operatorname{lk}(x)),
 \label{eq:B116-vertex-centralizer}\\
 C_{A(\Gamma)}(xy)&=\langle xy\rangle\times
 A(\operatorname{lk}(x)\cap\operatorname{lk}(y)).
 \label{eq:B116-pair-centralizer}
\end{align}
This proof uses the entire graph and ordinary HNN normal forms.

\emph{B.116.3. Return to the actual semidirect product.}
Write an element of $R$ as $u\delta$, with $u\in N$ and
$\delta=(g,z^m)\in D\times\mathbb Z$. The abelianization of $N$
is free abelian on all vertices of $\Omega$. Inner conjugation by
$u$ fixes that abelianization. Therefore $u\delta$ can centralize
$a$ only if $\delta$ fixes the coordinate $(T,a,0)$. This is precisely
$m=0$ and $g\in T$. Since the actual $k$ fixes $a$, we obtain
\[
 C_R(a)=C_N(a)\rtimes\langle k\rangle.
\]
The action fixes the vertex $b$ and shifts the integer $m$ in each
of the three other families in \eqref{eq:B116-link-a}. Eliminating
their orbit generators leaves exactly the presentation
\[
 \langle b,k,CaC,CdC,CeC\mid[k,b]=1\rangle
       =\mathbb Z^2(b,k)*F(CaC,CdC,CeC).
\]
The cyclic $a$ factor is fixed, which gives
\eqref{eq:B116-centralizer-a}.

For $aB$, abelianization has one $a$-type and one $b$-type coordinate.
The coarse action preserves the family type, so each of these two
coordinates must be fixed separately. It follows that $m=0$ and
$g\in T\cap CTC=1$. Thus every element centralizing $aB$ already
lies in $N$. Equations \eqref{eq:B116-common-link} and
\eqref{eq:B116-pair-centralizer} give
\eqref{eq:B116-centralizer-prefix}.

Their abelianizations are $\mathbb Z^6$ and $\mathbb Z^3$.
An automorphism mapping $a$ to $aB$ would restrict to an isomorphism
between these centralizers, which is impossible. The literal equality
$\theta(a)=aB$ was established using the actual neck marking in B.115.
No conclusion about an infinite-order centralizer in a completion is
needed or inferred.
\end{proof}
